\section{Discussion}
\label{sec:discussion}

The main comparison identifies an accuracy--storage operating point (Table~\ref{tab:main_results} and Fig.~\ref{fig:accuracy_storage}). On the 96 fixed-base conditions, EPOC reduces mean MSE by 15.40\% relative to Static with 6,352 median retained bytes. Full ELF increases that reduction to 19.29\%, but retains about 75 times as many bytes. EPOC also uses less state than the $\delta$-Adapter, COSA, FAC, and OMPB, and has lower paired MSE than each in a majority of conditions. The choice depends on the series and the available state budget: on BDG2 (Fox), the $\delta$-Adapter, COSA, and FAC have lower MSE with both backbones, while the width-2 and width-4 TEFL-style adapters in the separate joint-base experiment (Table~\ref{tab:learned_comparison}) retain fewer bytes than EPOC.

The endpoint controls clarify what the retained scalar contributes (Table~\ref{tab:endpoint_selection} and Fig.~\ref{fig:input_ablation}). At the same regression dimension and retained size, the preceding block's endpoint gives 1.65--2.20\% lower paired MSE than its first value, middle value, or mean. The low-order DCT coefficients summarize the completed block, whereas the observed endpoint explicitly records its latest residual; blocks with identical retained coefficients can have different endpoints. Without an endpoint input, the mean MSE reduction from Static falls to 13.37\%. Yet the endpoint projected from the retained coefficients achieves 14.98\%, close to the 15.40\% obtained with the observed endpoint. Much of the measured benefit therefore comes from supplying an endpoint-like feature to every component regression, with a smaller additional gain from retaining its uncompressed value in these experiments. The complete input also lowers paired MSE by 4.55\% relative to endpoint-only regression across the 96 conditions, supporting the combined use of residual coefficients, the shared endpoint feature, and current-forecast coefficients.

Among the tested settings, the number of retained components sets the principal accuracy--storage choice within EPOC (Fig.~\ref{fig:overview_sensitivity}). Increasing $K$ from 1 to 4 raises the mean MSE reduction from 11.13\% to 15.40\%, while increasing median retained state from 2,056 to 6,352 B. Increasing $K$ further to 8 adds 1.00 percentage point of MSE reduction and 5,728 B of state. Thus, $K=4$ is a compact reference setting rather than the most accurate tested setting; the preferred $K$ depends on the storage available for each channel.

The evidence covers eight series, two base forecasters, an input length of 96, a forecast horizon of 24, and updates after complete target blocks. Retained-byte counts describe persistent auxiliary numerical arrays between forecasts. Evaluating other horizons and series, and measuring peak memory, latency, and energy on target devices, would extend the accuracy--storage comparison to deployment costs.
